\documentclass{sener2025}

\title{Prueba de Auto-Fit Final}
\subtitle{Verificación de Tabularx y Xltabular}
\date{\today}
\setDocumentoCorto{AutoFitTest}

\begin{document}

\maketitle

\portadaseccion{1}{Prueba de Auto-Fit}{Tablas modernas}

\section{Tablas Cortas con Tabularx}

Esta tabla utiliza \texttt{tabularx} con el tipo de columna \texttt{XH} y \texttt{XG} para auto-ajuste.

\begin{tabladoradoCorto}
  \begin{tabularx}{\textwidth}{V v v}
    \toprule
    \rowcolor{gobmxDorado} \encabezadodorado{Concepto} & \encabezadodorado{Descripción Larga para Probar el Auto-Fit} & \encabezadodorado{Estatus} \\
    \midrule
    Proyecto A & Esta es una descripción bastante larga que debería ajustarse automáticamente al ancho de la columna sin desbordar el margen de la página. & \badge[gobmxVerde]{Activo} \\
    Proyecto B & Texto corto. & \badge[gobmxDorado]{Pendiente} \\
    \bottomrule
  \end{tabularx}
\end{tabladoradoCorto}

\section{Tablas Largas con Xltabular}

Esta tabla utiliza \texttt{xltabular} para permitir saltos de página y auto-ajuste simultáneo.

\begin{tabladoradoLargo}
  \begin{xltabular}{\textwidth}{Y Z Z}
    \caption{Prueba de tabla larga con xltabular}\label{tab:larga_x}\\
    \toprule
    \rowcolor{gobmxDorado} \encabezadodorado{Categoría} & \encabezadodorado{Detalle Técnico con Mucho Texto para Validar la Distribución} & \encabezadodorado{Valor} \\
    \midrule
    \endfirsthead
    
    \multicolumn{3}{l}{\textit{Continuación...}} \\
    \toprule
    \rowcolor{gobmxDorado} \encabezadodorado{Categoría} & \encabezadodorado{Detalle Técnico} & \encabezadodorado{Valor} \\
    \midrule
    \endhead
    
    Infraestructura & El despliegue de la red nacional de transmisión requiere una coordinación precisa entre las diferentes regiones de control del CENACE. & 95\% \\
    Generación & La integración de nuevas centrales de ciclo combinado ayudará a estabilizar el sistema durante los picos de demanda máxima. & 100\% \\
    Sustentabilidad & Se estima que la participación de energías limpias alcance la meta establecida para el cierre de este ciclo de planeación. & 31.2\% \\
    \bottomrule
  \end{xltabular}
\end{tabladoradoLargo}

\section{Verificación de Portada de Sección Vacía}

A continuación una portada de sección con argumentos vacíos para verificar que ya no falla.

\portadaseccion{2}{}{}

\section{Fin de Prueba}
La compilación llegó hasta aquí correctamente.

\end{document}
